\chapter{Energieverbrauchsberechnung}
\label{app:energieberechnung}

\section{Annahmen und Datengrundlage}

Die Berechnung des Energieverbrauchs basiert auf folgenden Annahmen:

\begin{itemize}
    \item \textbf{LLM-basierte Lösung}: 1,82 Wh pro API-Aufruf 
          (Durchschnittswert basierend auf \cite{strubellEnergyPolicyConsiderations2019} 
          und aktuellen Schätzungen für mittelgroße Sprachmodelle)
    \item \textbf{Interaktionsmuster}: 20 API-Aufrufe pro Formulardurchlauf 
          (2 Interaktionen pro Feld: Frage generieren + Antwort verarbeiten)
    \item \textbf{Klassisches NLP}: 20W CPU-Last für 0,021s pro Feld
    \item \textbf{Regelbasiertes System}: 15W CPU-Last für 0,005s pro Feld
\end{itemize}

\section{Einschränkungen}

Diese Berechnung dient als Größenordnungsabschätzung und unterliegt folgenden Limitationen:

\begin{itemize}
    \item Energieverbrauch von LLM-APIs ist herstellerabhängig und oft intransparent
    \item Netzwerklatenz und Übertragungskosten sind nicht eingerechnet
    \item Alternative Systeme könnten je nach Implementierung variieren
    \item Die Zahlen verstehen sich als konservative Schätzungen
\end{itemize}

Trotz dieser Unsicherheiten liegt die Größenordnung des Unterschieds bei Faktor $10^3$ bis $10^5$, was die grundsätzliche Aussage zur Ineffizienz bestätigt.